\documentclass[conference]{IEEEtran}

% --- Encoding & Language ---
\usepackage[utf8]{inputenc}
\usepackage[T1]{fontenc}
\usepackage[spanish]{babel}

% --- Tables & Figures ---
\usepackage{booktabs}
\usepackage{tabularx}
\usepackage{array}
\usepackage{float}

% --- Typography ---
\usepackage{enumitem}
\usepackage{hyperref}
\usepackage{url}
\usepackage{cite}

% --- Page numbers ---
\pagestyle{plain}

% --- Compact lists ---
\setlist{nosep, leftmargin=1.5em}

% ==========================================
% DOCUMENT
% ==========================================

\title{Respuesta a la contrapropuesta: Defensa y refinamiento del Modelo de Plataforma para la reutilizaci\'on de software}

\author{
  \IEEEauthorblockN{Gonzalo Enrique Guti\'errez Castillo}
  \IEEEauthorblockA{\'Area de Digitales -- ACIDE\\
  Arequipa, Per\'u\\
  gonzalog@acideperu.org\\
  gonzalogut.99@gmail.com}
}

\begin{document}
\maketitle

\begin{abstract}
Se presenta una respuesta formal a la contrapropuesta titulada
\textit{``Reutilizaci\'on basada en plantillas y composici\'on como
alternativa a librer\'ias y OOP''}. Se reconocen las cr\'iticas
leg\'itimas al dise\~no inicial, se corrigen imprecisiones
factuales sobre la implementaci\'on existente, y se demuestra
---~mediante fuentes autoritativas y evidencia de la industria~---
que las premisas centrales de la contrapropuesta contienen
malinterpretaciones de los principios YAGNI y KISS. Se concluye
que ambas propuestas convergen en el mismo mecanismo de
distribuci\'on (copiar y adaptar), y que la diferencia real es
terminol\'ogica, no arquitect\'onica.
\end{abstract}

\begin{IEEEkeywords}
reutilizaci\'on, arquitectura de software, YAGNI, KISS, Clean
Architecture, plataforma, plantillas, ERP.
\end{IEEEkeywords}

% ==========================================
\section{Introducci\'on}
% ==========================================

La contrapropuesta~\cite{araoz2026} presenta cr\'iticas al modelo
de reutilizaci\'on propuesto en~\cite{gutierrez2026}, planteando un
enfoque basado en plantillas y composici\'on como alternativa.
Este documento responde a dichas cr\'iticas con tres objetivos:

\begin{enumerate}
  \item Reconocer las cr\'iticas leg\'itimas y las fallas del
    primer informe.
  \item Demostrar, con fuentes autoritativas, que varias premisas
    de la contrapropuesta contienen malinterpretaciones de
    principios fundamentales de ingenier\'ia de software.
  \item Evidenciar que ambas propuestas convergen en la misma
    soluci\'on, y que la discusi\'on debiera enfocarse en los
    detalles de implementaci\'on, no en una falsa dicotom\'ia.
\end{enumerate}

% ==========================================
\section{Reconocimiento de cr\'iticas leg\'itimas}
% ==========================================

Es importante reconocer, con honestidad intelectual, los puntos
donde la contrapropuesta identifica problemas reales.

\subsection{El dise\~no de Persona/Estudiante fue insuficiente}

La cr\'itica al caso concreto de \texttt{Persona.Surname} versus
\texttt{Estudiante.ApellidoPaterno/ApellidoMaterno} es leg\'itima.
La abstracci\'on de Persona no contempl\'o adecuadamente la
realidad del dominio educativo peruano, donde el manejo de dos
apellidos separados es un requisito legal y administrativo.

Sin embargo, que una abstracci\'on est\'e mal dise\~nada no
invalida el concepto de abstracci\'on. El problema fue de
\textit{dise\~no}, no de \textit{principio}. La entidad unificada
de Persona/Contacto es un patr\'on consolidado en la industria del
software empresarial, adoptado independientemente por los
principales ERPs del mundo (ver Secci\'on~\ref{sec:erps}).

La soluci\'on correcta era dise\~nar \texttt{Persona} con campos
de nombre estructurados que soporten m\'ultiples formatos culturales
---~como lo hace SAP en su tabla \texttt{BUT000} con
\texttt{NAME\_FIRST}, \texttt{NAME\_LAST} y
\texttt{NAME\_LAST2}~--- en lugar de usar un \'unico campo
\texttt{Surname}.

\subsection{Falta de claridad operativa}

El primer informe present\'o el \textit{qu\'e} y el
\textit{por qu\'e}, pero no detall\'o suficientemente el
\textit{c\'omo}: gobernanza, ownership del c\'odigo, resoluci\'on
de conflictos entre librer\'ia y negocio, y procesos de
actualizaci\'on. Esto es un vac\'io leg\'itimo que este documento
no pretende minimizar.

\subsection{Composici\'on sobre herencia}

La recomendaci\'on de composici\'on sobre herencia es
t\'ecnicamente correcta y ampliamente aceptada~\cite{gof1994}.
No obstante, es necesario aclarar una imprecisi\'on factual:
\textbf{la implementaci\'on en Colver ya usaba composici\'on}.
\texttt{Estudiante} no hereda de \texttt{Persona}; hereda de
\texttt{AggregateRoot} (su ra\'iz de agregado en DDD) y mantiene
una referencia a \texttt{Persona} mediante \texttt{PersonId}
(clave for\'anea). La relaci\'on 1:1 entre \texttt{Estudiante} y
\texttt{Persona} es una relaci\'on de composici\'on, no de
herencia de clases.

\subsection{Sobre la iniciativa unilateral}

El punto sobre falta de consenso es v\'alido en su forma.
Sin embargo, es necesario contextualizar: la acci\'on fue tomada
porque, a pesar de que el equipo se compromet\'ia reiteradamente
con una visi\'on de reutilizaci\'on, no se materializaban acciones
concretas para avanzar hacia ese objetivo. Se manten\'ia el
\textit{status quo} proyecto tras proyecto. Ante la inacci\'on
colectiva, se opt\'o por presentar una implementaci\'on concreta
que sirviera como punto de partida para la discusi\'on ---~no
como imposici\'on, sino como catalizador.

% ==========================================
\section{Sobre YAGNI y KISS: alcance, uso correcto y malinterpretaciones}
\label{sec:yagni}
% ==========================================

Es necesario aclarar una posici\'on fundamental: \textbf{no nos
oponemos a YAGNI ni a KISS}. Son principios valiosos y necesarios
para el desarrollo de software. Sin embargo, no operan en
aislamiento: deben aplicarse \textbf{en armon\'ia} con DRY
(\textit{Don't Repeat Yourself}), SSOT (\textit{Single Source of
Truth}), los principios SOLID, el paradigma SPLE
(\textit{Software Product Line Engineering}) y SINE
(\textit{Simple Is Not Easy})~\cite{sine}. Este \'ultimo principio
es especialmente relevante: la simplicidad no es lo mismo que la
facilidad. Construir una soluci\'on simple requiere dise\~no
deliberado y esfuerzo; lo f\'acil es escribir c\'odigo r\'apido
sin estructura, pero eso no es simple, es deuda t\'ecnica.
Ning\'un principio tiene supremac\'ia absoluta sobre los dem\'as;
la ingenier\'ia de software consiste precisamente en encontrar el
equilibrio correcto seg\'un el contexto.

\subsection{D\'onde YAGNI y KISS s\'i aplican: el modelado de entidades}

Es importante ser precisos. La contrapropuesta no critica la
arquitectura (Clean Architecture, CQRS, DDD) ---~ambas partes
estamos de acuerdo en esos patrones~---. Lo que critica es el
\textbf{modelado de entidades}: que se haya modelado
\texttt{Branch} (m\'ultiples sedes) cuando Colver tiene una sola,
que \texttt{Persona} tenga un campo \texttt{Surname} gen\'erico
en lugar de apellidos espec\'ificos, y que se hayan dise\~nado
entidades pensando en m\'ultiples rubros durante el desarrollo de
un proyecto espec\'ifico.

En este nivel, YAGNI tiene legitimidad. Fowler define
\textit{``presumptive feature''} como \textit{``any code that
supports a feature that isn't yet being made available for
use''}~\cite{fowler2015yagni}. Si model\'aramos entidades para un
dominio que genuinamente desconocemos, eso ser\'ia especulaci\'on
y YAGNI aplicar\'ia con toda raz\'on.

Aceptamos esta cr\'itica parcialmente: durante el desarrollo de
Colver, el foco debi\'o estar en resolver Colver primero. Ah\'i
YAGNI y KISS son gu\'ias correctas.

\subsection{D\'onde YAGNI y KISS no aplican: reutilizaci\'on conocida}

Sin embargo, la contrapropuesta extiende YAGNI m\'as all\'a de su
alcance leg\'itimo cuando lo usa para rechazar \textit{toda}
inversi\'on en reutilizaci\'on. Jeffries, co-creador de XP y del
principio YAGNI, define su alcance as\'i:

\begin{quote}
\textit{``Always implement things when you \textbf{actually} need
them, never when you just \textbf{foresee} that you need
them.''}~\cite{jeffries_yagni}
\end{quote}

La palabra clave es \textit{actually} versus \textit{foresee}.
Cuando la consultora ya ejecut\'o m\'ultiples proyectos con los
mismos patrones (Persona, Auth, Auditor\'ia), la necesidad ya no
es especulativa: \textit{actually need} ya se cumpli\'o.

Fowler refuerza esta distinci\'on:

\begin{quote}
\textit{``If you do something for a future need that doesn't
actually increase the complexity of the software, then there's no
reason to invoke yagni.''}~\cite{fowler2015yagni}
\end{quote}

Modelar una entidad \texttt{Branch} con un FK no aumenta la
complejidad del sistema de forma significativa, pero habilita la
reutilizaci\'on para clientes con m\'ultiples sedes. Es importante
notar que respetar YAGNI aqu\'i no significa eliminar la entidad:
significa \textbf{no implementar l\'ogica de cambio de sede,
creaci\'on de nuevas sedes ni interfaces de gesti\'on multi-sede}
mientras el cliente no lo solicite. La entidad existe en el modelo
de datos; los features asociados no, hasta que se necesiten.

Este no es un escenario hipot\'etico. En un proyecto anterior de
la consultora ---~el sistema para Cl\'inica Juan Pablo~II
Dermatolog\'ia L\'aser y Est\'etica S.A.C.
(RUC~20603130295)~--- se desarroll\'o inicialmente sin soporte
multi-sede porque el cliente solicit\'o gesti\'on global. Cuando
posteriormente la cl\'inica requiri\'o operar por sedes, el
backend no pod\'ia refactorizarse sin riesgo, y fue necesario
aplicar parches inseguros que compromet\'ian la integridad del
sistema. Si la entidad \texttt{Branch} hubiera existido desde el
modelo de datos inicial ---~sin l\'ogica adicional, solo la
estructura~--- la adaptaci\'on habr\'ia sido un feature nuevo en
lugar de una cirug\'ia de emergencia.

La propuesta inicial de librer\'ias para Colver no surgi\'o de una
abstracci\'on arbitraria ni de especulaci\'on te\'orica: fue
construida sobre el conocimiento recogido en esta clase de
antecedentes reales. Los patrones que se propuso estandarizar
(Persona, Organizaci\'on, Sedes, Autenticaci\'on) son exactamente
los que generaron fricci\'on y parches en proyectos anteriores
cuando no estaban previstos en la estructura. Aplicar YAGNI a
este tipo de decisiones estructurales es confundir simplicidad con
recorte.

\subsection{YAGNI sin calidad es ``skimping''}

Jeffries distingue expl\'icitamente entre YAGNI leg\'itimo y
recorte de calidad:

\begin{quote}
\textit{``YAGNI, on its own, might not be safe\ldots\ But when we
are keeping the design quality, code quality, and test quality
high, YAGNI is quite safe.''}~\cite{jeffries2019}
\end{quote}

Su conclusi\'on es categ\'orica:

\begin{quote}
\textit{``YAGNI + clean code? Yes. Skimping?
No.''}~\cite{jeffries2019}
\end{quote}

Fowler complementa desde la perspectiva econ\'omica:

\begin{quote}
\textit{``Yagni is not a justification for neglecting the health
of your code base. Yagni requires (and enables) malleable
code.''}~\cite{fowler2015yagni}
\end{quote}

\subsection{Qu\'e significa aplicar YAGNI y KISS correctamente en la consultora}

Proponemos una interpretaci\'on razonable de ambos principios
para el contexto espec\'ifico de la consultora:

\begin{enumerate}
  \item \textbf{YAGNI para features del cliente}: No modelar
    funcionalidades que el cliente actual no ha solicitado. Si el
    colegio no necesita facturaci\'on electr\'onica hoy, no se
    construye. Esto es YAGNI leg\'itimo.

  \item \textbf{KISS para la implementaci\'on}: Cada feature se
    implementa de la forma m\'as simple que resuelva el problema.
    Sin capas innecesarias, sin abstracciones especulativas
    dentro del proyecto.

  \item \textbf{DRY y SSOT entre proyectos}: Cuando un patr\'on
    ya se repiti\'o en m\'ultiples proyectos (Persona, Auth,
    Auditor\'ia), se consolida en el repositorio base. Esto no
    viola YAGNI porque la necesidad ya est\'a demostrada; y no
    viola KISS porque reduce la complejidad global al eliminar
    duplicaci\'on.

  \item \textbf{SINE para el repositorio base}: Construir
    plantillas reutilizables bien dise\~nadas \textit{no es
    f\'acil}, pero el resultado \textit{s\'i es simple}. La
    simplicidad del producto final justifica el esfuerzo de
    dise\~no~\cite{sine}.
\end{enumerate}

Esta interpretaci\'on es coherente con las 4 reglas de dise\~no
simple de Kent Beck, ordenadas por
prioridad~\cite{fowler_beck}:

\begin{enumerate}
  \item Pasa los tests.
  \item Revela intenci\'on.
  \item \textbf{No tiene duplicaci\'on.}
  \item Tiene el m\'inimo de elementos.
\end{enumerate}

La regla 3 (no duplicaci\'on) tiene \textbf{mayor prioridad} que
la regla 4 (m\'inimos elementos). Copiar y pegar el mismo c\'odigo
de autenticaci\'on, personas y auditor\'ia en cada proyecto viola
la regla 3, que es m\'as importante que ``tener menos archivos''.
Bajo la definici\'on original de Beck, el repositorio base
\textbf{es m\'as simple} que duplicar c\'odigo entre proyectos.

\subsection{El retorno de la inversi\'on en dise\~no}

Fowler, en su \textit{Design Stamina Hypothesis}, demuestra que
la inversi\'on en dise\~no se recupera r\'apidamente:

\begin{quote}
\textit{``You can save short-term time by neglecting design, but
this accumulates Technical Debt which will slow your productivity
later.''}~\cite{fowler_stamina}
\end{quote}

\begin{quote}
\textit{``I take the view that [the design payoff line] is much
lower than most people think: usually weeks not
months.''}~\cite{fowler_stamina}
\end{quote}

Y sobre la calidad interna:

\begin{quote}
\textit{``High internal quality reduces the cost of future
features, meaning that putting the time into writing good code
actually reduces cost. [\ldots] The `cost' of high internal quality
software is negative.''}~\cite{fowler_quality}
\end{quote}

Robert C.\ Martin lo resume:

\begin{quote}
\textit{``If you think good architecture is expensive, try bad
architecture.''}~\cite{martin2017}
\end{quote}

% ==========================================
\section{Imprecisiones factuales en la contrapropuesta}
% ==========================================

\subsection{Se critica un strawman: librer\'ia r\'igida vs.\ modelo propuesto}

La contrapropuesta dedica la mayor parte de sus cr\'iticas al
concepto de ``librer\'ias r\'igidas que no se pueden tocar''. Sin
embargo, el informe original~\cite{gutierrez2026} explica
expl\'icitamente el modelo \textit{shadcn}:

\begin{quote}
\textit{``Cada proyecto toma lo que necesita y lo moldea
libremente. Las modificaciones en un proyecto no afectan a los
dem\'as.''}~\cite{gutierrez2026}
\end{quote}

La contrapropuesta propone exactamente lo mismo: plantillas que
se copian, se pegan y se adaptan. \textbf{Es el mismo mecanismo de
distribuci\'on con diferente terminolog\'ia.} La Tabla~\ref{tab:convergencia}
muestra la equivalencia.

\begin{table}[H]
\centering
\caption{Convergencia entre ambas propuestas}
\label{tab:convergencia}
\small
\begin{tabularx}{\columnwidth}{>{\raggedright\arraybackslash}X >{\raggedright\arraybackslash}X}
\toprule
\textbf{Propuesta original} & \textbf{Contrapropuesta} \\
\midrule
Repositorio base & Repo.\ de plantillas \\
Copiar y adaptar (shadcn) & Copiar y pegar \\
Clean Architecture + DDD & Clean Arch.\ + DDD \\
Proyecto due\~no del c\'odigo & Proyecto due\~no del c\'odigo \\
No modificar repo base por un cliente & Plantillas indep.\ de proyectos \\
\bottomrule
\end{tabularx}
\end{table}

\subsection{No es overengineering: los flujos eran conocidos}

La contrapropuesta argumenta que se intent\'o
\textit{``replicar completamente toda la estructura utilizada por
Odoo''}~\cite{araoz2026}. Esto es incorrecto. Las entidades
implementadas en las librer\'ias (Persona, Organizaci\'on,
Empleado, Autenticaci\'on) no son especulaciones te\'oricas: son
patrones que el equipo ya hab\'ia desarrollado en proyectos
anteriores.

George Fairbanks proporciona el marco adecuado para evaluar
cu\'ando la inversi\'on arquitect\'onica est\'a justificada:

\begin{quote}
\textit{``There is no need for meticulous designs when risks are
small, nor any excuse for sloppy designs when risks threaten your
success.''}~\cite{fairbanks2010}
\end{quote}

En el contexto de una consultora con m\'ultiples proyectos,
clientes repetidos y rotaci\'on de personal, el riesgo de no
estandarizar es alto. La inversi\'on arquitect\'onica est\'a
justificada.

\subsection{El ejemplo de ``Colver solo tiene una sede'' es falaz}

Que el primer cliente tenga una \'unica sede no invalida modelar
la entidad \texttt{Branch}. Esta entidad a\~nade un FK a la base
de datos y permite que el sistema funcione para clientes con $N$
sedes sin modificaci\'on alguna. El costo marginal es m\'inimo
(una tabla, una clave for\'anea) y el beneficio de
reutilizaci\'on es alto. Fowler lo expresa directamente:

\begin{quote}
\textit{``If you do something for a future need that doesn't
actually increase the complexity of the software, then there's no
reason to invoke yagni.''}~\cite{fowler2015yagni}
\end{quote}

% ==========================================
\section{La abstracci\'on de Persona: est\'andar de la industria}
\label{sec:erps}
% ==========================================

La contrapropuesta trata la entidad unificada de Persona como un
ejemplo de sobre-ingenier\'ia. Sin embargo, este patr\'on es el
est\'andar de facto en todo el software empresarial del mundo.
Cada uno de los principales ERPs lo ha adoptado de forma
independiente.

\subsection{Odoo: \texttt{res.partner}}

Odoo utiliza un \'unico modelo \texttt{res.partner} para
\textbf{todos} los contactos del sistema: individuos, empresas,
clientes, proveedores, empleados y direcciones. Los roles se
asignan mediante campos de rango: \texttt{customer\_rank} y
\texttt{supplier\_rank}~\cite{odoo_partner}. Cada m\'odulo del
ERP (Ventas, Compras, Contabilidad, Inventario, CRM, RRHH)
referencia \texttt{res.partner}. Una misma empresa puede ser
simult\'aneamente cliente y proveedor.

\subsection{SAP S/4HANA: Business Partner (\texttt{BUT000})}

SAP reemplaz\'o sus tablas maestras separadas
---~\texttt{KNA1} (cliente) y \texttt{LFA1}
(proveedor)~--- por un \textbf{Business Partner} unificado en
\texttt{BUT000}. Los roles se asignan en una tabla separada
\texttt{BUT100}, permitiendo $N$ roles por entidad con fechas de
validez~\cite{sap_bp}. Notablemente, \texttt{BUT000} soporta
expl\'icitamente m\'ultiples apellidos: \texttt{NAME\_LAST} y
\texttt{NAME\_LAST2}, dise\~nado para culturas
hispanoamericanas~\cite{sap_but000}.

\subsection{Microsoft Dynamics 365: \texttt{DirPartyTable}}

Microsoft documenta oficialmente el patr\'on:

\begin{quote}
\textit{``A party is a person or an organization that is involved
in a business. When you use the party concept, a person or an
organization can play more than one role in a business (for
example, customer, vendor, or contact).''}~\cite{ms_dynamics}
\end{quote}

Y justifica con un ejemplo concreto:

\begin{quote}
\textit{``A party can be both a customer and a vendor. [\ldots]
A party can be both an employee and a customer. For example, an
employee of Contoso buys electronics from Contoso for personal
use.''}~\cite{ms_dynamics}
\end{quote}

\subsection{Oracle TCA: \texttt{HZ\_PARTIES}}

Oracle implementa la \textit{Trading Community Architecture}
(TCA), quiz\'as la implementaci\'on m\'as rigurosa del patr\'on.
En TCA, un \texttt{Party} \textbf{no es} un cliente: se convierte
en cliente cuando se crea un registro en
\texttt{HZ\_CUST\_ACCOUNTS}. La misma entidad puede tener
m\'ultiples cuentas de cliente y ser simult\'aneamente proveedor.
Es un modelo de tres capas: Party (identidad), Account (rol) y
Site (ubicaci\'on)~\cite{oracle_tca}.

\subsection{ERPNext: Party Type din\'amico}

ERPNext utiliza el campo polim\'orfico \texttt{party\_type} +
\texttt{party} mediante \textit{Dynamic Links} del framework
Frappe. Direcciones y contactos se comparten entre Customer,
Supplier y Employee mediante este mecanismo~\cite{erpnext_party}.

\subsection{S\'intesis}

La Tabla~\ref{tab:erps} resume el patr\'on. \textbf{Todos} los
ERPs mayores del mundo usan una abstracci\'on unificada de
Persona/Entidad. Esto no es una coincidencia ni
sobre-ingenier\'ia: es la soluci\'on probada al problema real de
que una misma entidad del mundo real cumple m\'ultiples roles en
un sistema de negocio.

\begin{table}[H]
\centering
\caption{Patr\'on de Persona unificada en ERPs}
\label{tab:erps}
\small
\begin{tabularx}{\columnwidth}{l >{\raggedright\arraybackslash}X l}
\toprule
\textbf{ERP} & \textbf{Entidad central} & \textbf{Roles} \\
\midrule
Odoo & \texttt{res.partner} & Ranks \\
SAP & \texttt{BUT000} & \texttt{BUT100} \\
Dynamics & \texttt{DirPartyTable} & Party Roles \\
Oracle & \texttt{HZ\_PARTIES} & Accounts \\
ERPNext & \texttt{party\_type} & Dyn.\ Link \\
\bottomrule
\end{tabularx}
\end{table}

% ==========================================
\section{SPLE: el marco te\'orico que aplica}
% ==========================================

El paradigma de \textit{Software Product Line Engineering}
(SPLE), estandarizado por el Software Engineering Institute (SEI)
de Carnegie Mellon~\cite{clements2001}, define c\'omo construir un
portafolio de productos a partir de un n\'ucleo com\'un. La tesis
central es que cuando una organizaci\'on \textbf{sabe} que
construir\'a m\'ultiples productos similares, la inversi\'on
anticipada en arquitectura compartida genera retornos econ\'omicos
significativos.

La consultora opera exactamente en este escenario: m\'ultiples
proyectos, patrones comunes conocidos, y un objetivo estrat\'egico
expl\'icito de reutilizaci\'on. Esto contradice directamente la
aplicaci\'on ingenua de YAGNI en contextos multi-proyecto, porque
la reutilizaci\'on es \textbf{conocida y planificada}, no
especulativa.

% ==========================================
\section{An\'alisis de la analog\'ia Taller vs.\ F\'abrica}
% ==========================================

La contrapropuesta argumenta que la analog\'ia de taller versus
f\'abrica es inadecuada porque ``no fabricamos un \'unico modelo
de auto''~\cite{araoz2026}. La analog\'ia se extiende a fabricar
chasis, ejes y llantas, concluyendo que la f\'abrica
``acabar\'ia fabricando partes individuales peque\~nas'', y que
esto ser\'ia un desatino.

Sin embargo, eso es \textbf{exactamente} lo que hace la industria
automotriz real. En este contexto, una \textit{plataforma} es el
conjunto de componentes base que comparten varios modelos de auto:
la estructura del piso, las suspensiones, la ubicaci\'on del motor,
y los sistemas el\'ectricos fundamentales. No es un auto completo:
es la base com\'un sobre la que se construyen autos muy diferentes.

La Tabla~\ref{tab:autos} muestra ejemplos concretos. Un
Volkswagen Polo (hatchback econ\'omico) y un Audi Q3 (SUV
premium) comparten la plataforma MQB. Un Chevrolet Onix
(sed\'an), un Tracker (SUV) y una Montana (pickup) comparten la
plataforma GEM. Un Toyota Corolla y un RAV4 comparten TNGA. Un
Hyundai Tucson y un Kia Sportage ---~de marcas distintas~---
comparten la plataforma N3. Un Renault Duster y un Nissan Kicks
comparten CMF-B~\cite{vw_mqb, gm_gem, toyota_tnga,
hyundai_platform, renault_cmf}.

Nadie dir\'ia que un Toyota Corolla y un RAV4 son ``el mismo
auto'', ni que un Hyundai Tucson y un Kia Sportage son ``el mismo
veh\'iculo''. Son productos muy diferentes, para clientes muy
diferentes, a precios muy diferentes. Pero comparten una base com\'un que les permite a
las automotrices no reinventar lo fundamental en cada modelo nuevo.

La analog\'ia con la consultora es directa: un sistema para un
colegio y un sistema para un hotel son productos diferentes, pero
ambos necesitan autenticaci\'on, gesti\'on de personas y
auditor\'ia. Esos componentes son la ``plataforma'' de la
consultora. No es un sistema completo para todos: es la base
probada sobre la que cada proyecto construye su funcionalidad
espec\'ifica.

\begin{table}[H]
\centering
\caption{Plataformas compartidas en la industria automotriz}
\label{tab:autos}
\small
\begin{tabularx}{\columnwidth}{l >{\raggedright\arraybackslash}X r}
\toprule
\textbf{Plataforma} & \textbf{Modelos en LATAM} & \textbf{Ahorro} \\
\midrule
VW MQB & Polo, T-Cross, Tiguan, Taos, Audi A3/Q3 & 30\% \\
GM GEM & Onix, Tracker, Montana & -- \\
Toyota TNGA & Corolla, Corolla Cross, RAV4, Camry & 20\% \\
Hyundai N3/K2 & Tucson, Sportage, Creta, Seltos & -- \\
CMF-B (RNM) & Duster, Kicks, Captur & 30--40\% \\
\bottomrule
\end{tabularx}
\end{table}

% ==========================================
\section{Sobre la restricci\'on de uso de IA}
% ==========================================

La contrapropuesta establece que ``la idea original, el flujo de
trabajo, las entidades, sus nombres, propiedades y operaciones no
pueden venir de una IA''~\cite{araoz2026}. Si bien la intenci\'on
de evitar abstracciones especulativas es razonable, el criterio
es arbitrario.

La calidad de una abstracci\'on se mide por su resultado
(cohesi\'on, acoplamiento, expresividad del dominio), no por su
origen. Un humano puede producir abstracciones especulativas
tan deficientes como una IA, y una IA con el contexto adecuado
puede identificar patrones que un humano pasa por alto. El
criterio correcto es \textbf{revisi\'on y validaci\'on por el
equipo}, independientemente de qui\'en o qu\'e gener\'o la
propuesta inicial.

% ==========================================
\section{Sobre las limitaciones reconocidas}
% ==========================================

La contrapropuesta reconoce honestamente sus propias
limitaciones~\cite{araoz2026}: normalizaci\'on con composici\'on,
dependencias opcionales entre plantillas, infraestructura
(tablas propias vs.\ datos embebidos), manejo de errores, y
aspectos de DDD como Aggregates, Bounded Contexts y Eventos.

Es notable que estos son \textbf{exactamente los problemas que un
repositorio base institucional resuelve} al definir convenciones
claras desde el inicio. La contrapropuesta identifica vac\'ios
en su propia propuesta que la propuesta original ya aborda.

% ==========================================
\section{Conclusiones}
% ==========================================

\begin{enumerate}
  \item Se reconocen las cr\'iticas leg\'itimas al primer
    informe: el dise\~no insuficiente de Persona, la falta de
    detalle operativo, y la necesidad de mayor consenso.

  \item Se demuestra, con citas textuales de los autores
    originales (Fowler, Jeffries, Beck, Martin), que YAGNI no
    aplica a decisiones arquitect\'onicas ni a esfuerzos de
    mantenibilidad. Usar YAGNI para rechazar inversi\'on
    arquitect\'onica es una malinterpretaci\'on del principio.

  \item Se evidencia que la abstracci\'on unificada de
    Persona/Contacto es el est\'andar de facto en la industria,
    adoptado independientemente por Odoo, SAP, Microsoft, Oracle
    y ERPNext.

  \item Se demuestra que ambas propuestas convergen en el mismo
    mecanismo de distribuci\'on (copiar, adaptar, ser due\~no del
    c\'odigo), y que la diferencia es terminol\'ogica.

  \item La industria automotriz valida el modelo de plataforma:
    Toyota, Volkswagen, GM, Hyundai-Kia y Renault-Nissan
    construyen veh\'iculos radicalmente distintos sobre bases
    compartidas, reduciendo costos entre un 20\% y 40\%.

  \item Se propone que la discusi\'on futura se enfoque en los
    detalles de implementaci\'on (gobernanza, contenido de las
    plantillas, proceso de revisi\'on) en lugar de perpetuar una
    falsa dicotom\'ia entre enfoques que son, en esencia, el
    mismo.
\end{enumerate}

% ==========================================
% REFERENCES
% ==========================================

\begin{thebibliography}{99}

\bibitem{araoz2026}
F.~E.~Araoz~Morales, ``Reutilizaci\'on basada en plantillas y
composici\'on como alternativa a librer\'ias y OOP para la
reutilizaci\'on de c\'odigo,'' \'Area de Digitales -- ACIDE,
Arequipa, Per\'u, 2026.

\bibitem{gutierrez2026}
G.~E.~Guti\'errez~Castillo, ``Reporte T\'ecnico y
Estrat\'egico: Est\'andar Arquitect\'onico para la
Reutilizaci\'on de Software en una Consultora,'' \'Area de
Digitales -- ACIDE, Arequipa, Per\'u, 2026.

\bibitem{fowler2015yagni}
M.~Fowler, ``Yagni,'' \textit{martinfowler.com}, 2015.
[Online]. Available:
\url{https://martinfowler.com/bliki/Yagni.html}. [Accessed:
16-Mar-2026].

\bibitem{jeffries2019}
R.~Jeffries, ``YAGNI, Skimping, and Technical Debt,''
\textit{ronjeffries.com}, 2019. [Online]. Available:
\url{https://ronjeffries.com/articles/019-01ff/iter-yagni-skimp/}.
[Accessed: 16-Mar-2026].

\bibitem{jeffries_yagni}
R.~Jeffries, ``You're NOT Gonna Need It!,''
\textit{ronjeffries.com}. [Online]. Available:
\url{https://ronjeffries.com/xprog/articles/practices/pracnotneed/}.
[Accessed: 16-Mar-2026].

\bibitem{martin2017}
R.~C.~Martin, \textit{Clean Architecture: A Craftsman's Guide to
Software Structure and Design}.\hskip 1em plus 0.5em minus
0.4em\relax Pearson, 2017.

\bibitem{fowler_stamina}
M.~Fowler, ``Design Stamina Hypothesis,''
\textit{martinfowler.com}. [Online]. Available:
\url{https://martinfowler.com/bliki/DesignStaminaHypothesis.html}.
[Accessed: 16-Mar-2026].

\bibitem{fowler_quality}
M.~Fowler, ``Is High Quality Software Worth the Cost?,''
\textit{martinfowler.com}. [Online]. Available:
\url{https://martinfowler.com/articles/is-quality-worth-cost.html}.
[Accessed: 16-Mar-2026].

\bibitem{fowler_beck}
M.~Fowler, ``BeckDesignRules,'' \textit{martinfowler.com}.
[Online]. Available:
\url{https://martinfowler.com/bliki/BeckDesignRules.html}.
[Accessed: 16-Mar-2026].

\bibitem{fairbanks2010}
G.~Fairbanks, \textit{Just Enough Software Architecture: A
Risk-Driven Approach}.\hskip 1em plus 0.5em minus 0.4em\relax
Marshall \& Brainerd, 2010.

\bibitem{clements2001}
P.~Clements and L.~Northrop, \textit{Software Product Lines:
Practices and Patterns}.\hskip 1em plus 0.5em minus 0.4em\relax
Addison-Wesley, 2001.

\bibitem{gof1994}
E.~Gamma, R.~Helm, R.~Johnson, and J.~Vlissides, \textit{Design
Patterns: Elements of Reusable Object-Oriented Software}.\hskip
1em plus 0.5em minus 0.4em\relax Addison-Wesley, 1994.

\bibitem{odoo_partner}
Odoo S.A., ``res.partner source code,'' \textit{GitHub:
odoo/odoo}, branch 17.0. [Online]. Available:
\url{https://github.com/odoo/odoo/blob/17.0/odoo/addons/base/models/res_partner.py}.
[Accessed: 16-Mar-2026].

\bibitem{sap_bp}
SAP SE, ``Business Partner in SAP S/4HANA,'' \textit{SAP Help
Portal}. [Online]. Available:
\url{https://help.sap.com/docs/SAP_S4HANA_ON-PREMISE}.
[Accessed: 16-Mar-2026].

\bibitem{sap_but000}
``SAP BUT000 Table -- Business Partner: General Data,''
\textit{leanx.eu}. [Online]. Available:
\url{https://leanx.eu/en/sap/table/but000.html}. [Accessed:
16-Mar-2026].

\bibitem{ms_dynamics}
Microsoft, ``Party and global address book,'' \textit{Microsoft
Learn}. [Online]. Available:
\url{https://learn.microsoft.com/en-us/dynamics365/fin-ops-core/dev-itpro/data-entities/dual-write/party-gab}.
[Accessed: 16-Mar-2026].

\bibitem{oracle_tca}
Oracle, ``Trading Community Architecture (TCA),'' \textit{Oracle
Documentation}. [Online]. Available:
\url{https://docs.oracle.com/en/applications/}. [Accessed:
16-Mar-2026].

\bibitem{erpnext_party}
Frappe Technologies, ``ERPNext party.py source code,''
\textit{GitHub: frappe/erpnext}. [Online]. Available:
\url{https://github.com/frappe/erpnext/blob/develop/erpnext/accounts/party.py}.
[Accessed: 16-Mar-2026].

\bibitem{vw_mqb}
Volkswagen AG, ``Modular toolkit strategy as recipe for success:
the MQB celebrates tenth anniversary,'' \textit{Volkswagen
Newsroom}, 2022. [Online]. Available:
\url{https://www.volkswagen-newsroom.com/en/press-releases/modular-toolkit-strategy-as-recipe-for-success-the-mqb-celebrates-tenth-anniversary-8030}.
[Accessed: 16-Mar-2026].

\bibitem{gm_gem}
``GM GEM Vehicle Platform,'' \textit{GM Authority}. [Online].
Available: \url{https://gmauthority.com/blog/gm/gm-platforms/gem/}.
[Accessed: 16-Mar-2026].

\bibitem{toyota_tnga}
Toyota Motor Corporation, ``Toyota New Global Architecture,''
\textit{global.toyota}. [Online]. Available:
\url{https://global.toyota/en/mobility/tnga/index.html}.
[Accessed: 16-Mar-2026].

\bibitem{hyundai_platform}
``Hyundai-Kia N platform,'' \textit{Wikipedia}. [Online].
Available:
\url{https://en.wikipedia.org/wiki/Hyundai-Kia_N_platforms}.
[Accessed: 16-Mar-2026].

\bibitem{renault_cmf}
Renault-Nissan-Mitsubishi Alliance, ``Common Module Family
(CMF),'' 2013. [Online]. Available:
\url{https://en.wikipedia.org/wiki/Renault\%E2\%80\%93Nissan_Common_Module_Family}.
[Accessed: 16-Mar-2026].

\bibitem{sine}
M.~Lehtinen, ``4 Most Important Software Development Principles:
DRY, YAGNI, KISS, and SINE,'' \textit{mattilehtinen.com}.
[Online]. Available:
\url{https://mattilehtinen.com/articles/4-most-important-software-development-principles-dry-yagni-kiss-and-sine/}.
[Accessed: 16-Mar-2026].

\end{thebibliography}

\end{document}
